\documentclass[12pt]{article}
\usepackage[utf8]{inputenc}
\usepackage[T1]{fontenc}
\usepackage{hyperref}
\title{Wanukaka Language Resources for K-12 Curriculum Development}
\author{Compiled by Claude Sonnet 5.5 (claude-sonnet-5-5) for Andr\'as Kornai}
\date{}
\begin{document}
\maketitle

\section{Introduction}
Wanukaka (ISO 639-3: wnk; Glottocode wanu1241) is an Austronesian language of Sumba, East Nusa Tenggara, Indonesia. It is spoken mainly in the Wanokaka district of West Sumba. Dialects are given as Rua and Wanukaka. It is closest to Anakalang. Estimates put it at about 10,000 speakers. Searches found no school materials in Wanukaka. The language is not known to be taught in schools, and the sources found are encyclopedic, lexical and academic.

\section{Primary Educational Resources}
No primers, textbooks or school programmes in Wanukaka turned up in the searches. A curriculum would therefore have to be built from the lexical and linguistic sources below.

\section{Linguistic and Reference Materials}
\begin{itemize}
\item Wanukaka dictionary (Dictionaria): \url{https://dictionaria.clld.org/contributions/wanukaka}
\item Wikipedia overview of the Wanukaka language: \url{https://en.wikipedia.org/wiki/Wanukaka_language}
\item Phonological evidence distinguishing Wanokaka from Anakalang (Sumba Tengah, NTT; journal article in Indonesian): \url{https://ejournal.undiksha.ac.id/index.php/PRASI/article/view/6826}
\item Related record in the Indonesian GARUDA index: \url{https://garuda.kemdiktisaintek.go.id/documents/detail/6082956}
\item Verbix page for Wanukaka: \url{https://docs.verbix.com/Languages/Wanukaka}
\end{itemize}

\section{Religious and Community Resources}
Joshua Project has a people-group profile for the Wanukaka of Indonesia: \url{https://joshuaproject.net/people_groups/15859}. No scripture translations or community organisations were found in the searches.

\section{Digital and Audio Resources}
No audio, video or corpus resources were found in the searches.

\section{Curriculum Development Context}
Direct evidence of how the language is used is lacking. It is thought to be a first language of the whole ethnic community. It is not known to be taught in schools. Whether Wanukaka is a language or a dialect of a wider Sumbanese language is debated in the literature. Any curriculum work would need new documentation with local speakers, starting from the Dictionaria wordlist.
\end{document}